
\documentclass[11pt,a4paper]{article}
\usepackage[utf8]{inputenc}
\usepackage[brazil]{babel}
\usepackage{geometry}
\usepackage{booktabs}
\usepackage{array}
\usepackage{longtable}
\usepackage{enumitem}
\usepackage{hyperref}
\usepackage{xurl}
\urlstyle{tt}
\usepackage{amsmath}
\usepackage{tikz}
\usetikzlibrary{shapes.geometric,arrows.meta,positioning,fit,calc}

\geometry{margin=2.5cm}

\title{Manual Didático e Argumentativo\\do Modelo de Otimização de Mix}
\author{Mantiqueira -- 7D Analytics}
\date{\today}

\begin{document}

\maketitle
\tableofcontents
\newpage

\section{Introdução: por que este modelo existe}

Toda decisão de mix de produção enfrenta um dilema permanente: a empresa precisa
maximizar margem, mas não pode ignorar restrições físicas, compromissos
comerciais e a realidade da demanda. Em outras palavras, a melhor decisão
financeira no papel pode ser impraticável na operação; por outro lado, uma
decisão puramente operacional pode preservar rotina, mas destruir valor.

Este modelo nasce para resolver esse conflito de forma disciplinada. Seu papel
não é apenas calcular números, mas transformar uma escolha normalmente subjetiva
em um processo racional, auditável e repetível. Ao explicitar regras, limites e
efeitos, o modelo cria um padrão decisório que protege a margem sem romper o
atendimento e sem mascarar incertezas dos dados.

Este manual foi escrito com esse mesmo princípio: apresentar o modelo de forma
didática, porém argumentativa. Em vez de listar apenas ``o que fazer'',
explicamos também \textbf{por que cada decisão existe}, \textbf{o que ela busca
alcançar} e \textbf{qual risco ela evita}.

\section{Tese central do método}

A tese central do modelo é simples e forte:
\begin{quote}
\textbf{A otimização só gera valor de forma sustentável quando combina quatro
pilares simultaneamente: eficiência econômica, viabilidade operacional,
coerência comercial e transparência analítica.}
\end{quote}

Se qualquer pilar for ignorado, surgem distorções:
\begin{itemize}[leftmargin=1.2cm]
    \item sem eficiência econômica, há produção mal distribuída e margem perdida;
    \item sem viabilidade operacional, aparecem recomendações inexequíveis;
    \item sem coerência comercial, aumenta o risco de alocar volume sem saída;
    \item sem transparência, o resultado pode até estar correto, mas não é confiável para gestão.
\end{itemize}

\section{Problema de negócio em linguagem prática}

A capacidade produtiva é dada por classe. A pergunta gerencial não é ``quanto
produzir'', mas ``\textit{como distribuir o que foi produzido}''. Esse detalhe
é decisivo: como os SKUs de uma mesma classe têm margens e demandas diferentes,
a distribuição do volume altera diretamente o resultado financeiro.

Sem método, a distribuição tende a seguir hábito, pressão de curto prazo ou
repetição histórica sem crítica. Esse padrão geralmente:
\begin{itemize}[leftmargin=1.2cm]
    \item subutiliza SKUs mais rentáveis;
    \item superaloca SKUs com baixa tração;
    \item mistura reserva de pedido com alocação discricionária;
    \item dificulta explicar por que um item entrou e outro ficou fora.
\end{itemize}

O modelo substitui essa dinâmica por uma decisão com critérios explícitos e
mensuráveis.

\section{Guia rapido (5 minutos)}
\label{sec:guia-rapido}

Para uso diario de operacao/comercial, este e o roteiro minimo:
\begin{enumerate}[leftmargin=1.2cm]
    \item revisar \texttt{config.yaml} (periodo, estabelecimento, granularidade, pedidos; detalhe na Secao~\ref{sec:param-config});
    \item escolher o fluxo de execucao: \textbf{A} (completo), \textbf{B} (sem preparacao) ou \textbf{C} (parcial). O que cada um faz esta no diagrama da Secao~\ref{sec:casos-execucao}; os comandos exatos estao na Secao~\ref{sec:como-executar};
    \item rodar pipeline e abrir a pasta da rodada em \texttt{resultados/} (padrao de nome na Secao~\ref{sec:rastreabilidade});
    \item validar as tres saidas principais listadas na Secao~\ref{sec:leitura-resultados}: \texttt{comparacao\_producao\_alocacao}, \texttt{auditoria\_baseline}, \texttt{skus\_fora\_otimizacao};
    \item decidir go/no-go com o checklist da Secao~\ref{sec:checklist-executivo}.
\end{enumerate}

\textbf{Nao fazer (erros comuns):}
\begin{itemize}[leftmargin=1.2cm]
    \item alterar \texttt{data\_ref}/\texttt{semana\_ref} e rodar fluxo B sem garantir que os CSVs em \texttt{inputs/} estao no mesmo periodo;
    \item trocar \texttt{estabelecimentos} sem revisar pedidos/producao para o mesmo escopo;
    \item interpretar diferenca de alocacao sem separar reserva, producao liquida e itens fora.
\end{itemize}

\textbf{Glossario operacional curto:}
\begin{itemize}[leftmargin=1.2cm]
    \item \textbf{Baseline:} cenario de referencia (sem decisao otimizada), usado para medir ganho;
    \item \textbf{Comparador:} tabela de confronto producao vs alocacao por SKU;
    \item \textbf{PBI consolidado:} pacote final para BI, incluindo aba de SKUs fora;
    \item \textbf{DOW:} distribuicao historica por dia da semana (apenas granularidade semanal).
\end{itemize}

\section{Conceitos fundamentais (autocontido)}

\begin{longtable}{>{\raggedright\arraybackslash}p{5.8cm}>{\raggedright\arraybackslash}p{9.2cm}}
\toprule
\textbf{Conceito} & \textbf{Definição operacional} \\
\midrule
\endhead
SKU & Item comercial considerado na decisão (produto/embalagem). \\
Classe & Grupo produtivo no qual a capacidade é controlada. \\
Pedido garantido & Pedido confirmado, tratado como prioridade de atendimento. \\
Reserva & Volume separado para pedido antes da otimização do excedente. \\
Excedente & Volume restante da classe após reserva; é o volume elegível à otimização. \\
\url{atender_pedidos} & Parâmetro que liga/desliga a lógica de reserva prévia de pedidos. \\
\url{capar_reserva_na_producao} & Parâmetro que impede reserva acima da produção da classe no período. \\
\url{quantidade_alocada} & Volume definido pelo otimizador para o SKU. \\
\url{quantidade_reservada} & Volume do SKU atendido como compromisso de pedido. \\
\url{quantidade_nao_atendida_pedido} & Parcela do pedido que não pôde ser atendida no período. \\
Produção bruta / estorno EAC / produção líquida & Bruta = produzido; estorno EAC = cancelamento; líquida = bruta - estorno. \\
Fallback & Regra de preenchimento para dados ausentes (preço, custo, embalagem). \\
\bottomrule
\end{longtable}

\section{Arquitetura de decisão: do dado à ação}

\subsection{Etapa 1 -- Preparação (ETL)}

O ETL consolida os insumos que sustentam a decisão: produção por classe, custos,
preços, pedidos, demanda histórica, cadastro de classes e compatibilidade de
embalagem. Aqui, a principal discussão não é performance, mas confiabilidade.

Argumento central desta etapa: \textbf{não existe boa otimização com base
mal preparada}. Por isso, o pipeline:
\begin{itemize}[leftmargin=1.2cm]
    \item padroniza unidades;
    \item aplica filtros de elegibilidade;
    \item trata estorno (EAC) para calcular produção líquida;
    \item registra origem de preço, custo e embalagem para auditoria.
\end{itemize}

\subsection{Etapa 2 -- Otimização}

Com os dados prontos, o modelo distribui o excedente por SKU maximizando margem,
sob restrições de capacidade e, quando configurado, de demanda histórica.

A decisão não é ``livre'': ela é deliberadamente condicionada para impedir
soluções financeiramente atraentes, mas comercialmente inviáveis.

\subsection{Etapa 3 -- Explicação e governança}

Após resolver, o processo produz relatórios de comparação, auditoria e saída para
BI. Essa etapa não é acessória. Ela existe porque decisão sem explicação gera
resistência, retrabalho e baixa adoção.

\section{Fluxogramas: componentes, conversas e casos de execução}

Os diagramas abaixo refletem o encadeamento real do repositório: o shell
\texttt{executar\_pipeline.sh} dispara a preparação opcional dos CSVs em
\texttt{inputs/} e sempre encerra com \texttt{main.py}, que lê
\texttt{config.yaml}, executa ETL, otimização, comparativos e grava tudo em
\texttt{resultados/<timestamp>\_<D|S|M>/}. Cada script Python de preparação também lê o
mesmo \texttt{config.yaml}, garantindo período, estabelecimento e caminhos
alinhados.

\subsection{Visão macro (do parâmetro ao artefato)}

\begin{figure}[htbp]
\centering
\begin{tikzpicture}[
  node distance=6mm and 10mm,
  box/.style={rectangle, draw, rounded corners=2pt, minimum width=3.0cm,
    minimum height=9mm, align=center, font=\small},
  io/.style={rectangle, draw, dashed, rounded corners=2pt, minimum width=2.8cm,
    align=center, font=\footnotesize},
  decision/.style={diamond, draw, aspect=2.4, inner sep=2pt, align=center, font=\footnotesize},
  lbl/.style={font=\scriptsize, inner sep=1pt},
  arrow/.style={-{Stealth[length=2.2mm]}, thick}
]
  \node[io] (cfg) {\textbf{config.yaml}\\(dados, modelo, paths)};
  \node[box, below=of cfg] (venv) {Ambiente\\\texttt{venv} + \texttt{python3}};
  \node[box, below=of venv] (sh) {\texttt{executar\_pipeline.sh}};
  \node[decision, below=of sh] (dprep) {Flag\\\texttt{{-}{-}sem-preparacao}?};
  \node[box, below left=12mm and 4mm of dprep] (prep) {Preparação\\\textbf{passos 1--6}\\(gera/atualiza\\\texttt{inputs/*.csv})};
  \node[box, below right=12mm and 4mm of dprep] (skip) {Pula passos 1--6\\(mantém CSVs\\atuais)};
  \node[box, below=22mm of dprep] (main) {\textbf{main.py}\\Fases 1--6\\(ver detalhe)};
  \node[io, below=of main] (out) {\textbf{resultados/}\\\texttt{<timestamp>/}\\+ \texttt{pbi\_concat/}};
  \draw[arrow] (cfg) -- (venv);
  \draw[arrow] (venv) -- (sh);
  \draw[arrow] (sh) -- (dprep);
  \draw[arrow] (dprep) -- node[lbl, left, pos=0.35] {Não} (prep);
  \draw[arrow] (dprep) -- node[lbl, right, pos=0.35] {Sim} (skip);
  \draw[arrow] (prep) |- ([yshift=-2mm]main.west);
  \draw[arrow] (skip) |- ([yshift=-2mm]main.east);
  \draw[arrow] (main) -- (out);
\end{tikzpicture}
\caption{Visão macro: o \texttt{config.yaml} parametriza toda a cadeia; o shell
decide se os insumos em \texttt{inputs/} são regenerados antes do
\texttt{main.py}.}
\label{fig:fluxo-macro}
\end{figure}

\subsection{Detalhe --- preparação de insumos (passos 1--6)}

Cada etapa consome bases brutas indicadas em \texttt{paths.*} e grava um CSV
(ou planilha intermediária) sob \texttt{inputs/}, que o ETL do \texttt{main.py}
passa a tratar como contrato de entrada.

\begin{figure}[htbp]
\centering
\begin{tikzpicture}[
  node distance=5mm,
  prepstep/.style={rectangle, draw, rounded corners=2pt, minimum width=6.8cm,
    minimum height=8mm, align=left, font=\small},
  arrow/.style={-{Stealth[length=2mm]}, thick}
]
  \node[prepstep] (s1) {\textbf{1.} \texttt{extrair\_compatibilidade\_embalagem.py}
    $\rightarrow$ \texttt{compatibilidade\_sku\_embalagem.csv}};
  \node[prepstep, below=of s1] (s2) {\textbf{2.} \texttt{extrair\_precos\_embalagem.py}
    $\rightarrow$ \texttt{precos\_sku\_embalagem.csv}};
  \node[prepstep, below=of s2] (s3) {\textbf{3.} \texttt{gerar\_custos\_sku.py}
    $\rightarrow$ \texttt{custos\_sku.csv}};
  \node[prepstep, below=of s3] (s4) {\textbf{4.} \texttt{gerar\_demanda\_historica.py}
    $\rightarrow$ \texttt{demanda\_historica.csv}};
  \node[prepstep, below=of s4] (s5) {\textbf{5.} \texttt{gerar\_pedidos\_clientes.py}
    $\rightarrow$ \texttt{pedidos\_clientes.csv}};
  \node[prepstep, below=of s5] (s6) {\textbf{6.} \texttt{gerar\_producao\_classe.py}
    $\rightarrow$ \texttt{producao\_classe.csv}};
  \node[prepstep, below=of s6, draw=none, font=\itshape\small] (note)
    {Demais cadastros (classes, produção bruta, faturamento, etc.)\ permanecem
    referenciados em \texttt{paths} e são lidos no ETL.};
  \foreach \a/\b in {s1/s2,s2/s3,s3/s4,s4/s5,s5/s6,s6/note} {
    \draw[arrow] (\a) -- (\b);
  }
\end{tikzpicture}
\caption{Cadeia fixa de preparação: ordem importa quando um passo depende de
artefato anterior ou de mesma janela temporal no \texttt{config.yaml}.}
\label{fig:fluxo-preparacao}
\end{figure}

\subsection{Detalhe --- \texttt{main.py}: fases internas e dependências}

O orquestrador Python troca dados por meio de objetos e arquivos: o ETL devolve
\texttt{base\_otimizacao} e séries auxiliares (produção, excedente, pedidos); a
base é persistida em disco para auditoria; o otimizador consome esses insumos e
devolve alocação; em seguida entram pós-processamentos, comparativos e saídas.
A comparação produção\,×\,alocação é invocada por \texttt{subprocess}; o
consolidado PBI só roda se essa etapa tiver sucesso.

\begin{figure}[htbp]
\centering
\resizebox{\textwidth}{!}{%
\begin{tikzpicture}[
  node distance=4mm and 8mm,
  ph/.style={rectangle, draw, rounded corners=2pt, minimum height=9mm,
    align=center, font=\footnotesize, inner sep=3pt},
  arrow/.style={-{Stealth[length=1.8mm]}, thick}
]
  \node[ph] (c) {Carrega\\\texttt{config.yaml}};
  \node[ph, right=of c] (dir) {Cria\\\texttt{resultados/<ts>/}};
  \node[ph, right=of dir] (etl) {\textbf{Fase 1}\\ETLPipeline};
  \node[ph, right=of etl] (base) {Persiste\\\texttt{base\_otimizacao}\\(.parquet/.csv/.xlsx)};
  \node[ph, below=8mm of etl] (opt) {\textbf{Fase 2}\\Otimizador\\(OR-Tools)};
  \node[ph, right=of opt] (pos) {Pós:\\reserva + OUTROS};
  \node[ph, right=of pos] (bl) {\textbf{Fase 3}\\Baseline};
  \node[ph, below=6mm of pos] (sv) {\textbf{Fase 4}\\\texttt{salvar\_resultados}};
  \node[ph, left=of sv] (dow) {DOW\\(se gran.\ S)};
  \node[ph, right=of sv] (cmp) {\textbf{Fase 5}\\\texttt{comparar\_producao\_alocacao}};
  \node[ph, below=6mm of sv] (pbi) {\textbf{Fase 6}\\\texttt{gerar\_consolidado\_pbi}\\(se Fase 5 OK)};
  \draw[arrow] (c) -- (dir);
  \draw[arrow] (dir) -- (etl);
  \draw[arrow] (etl) -- (base);
  \draw[arrow] (etl) -- node[right, font=\scriptsize] {\texttt{df\_base}} (opt);
  \draw[arrow] (opt) -- (pos);
  \draw[arrow] (pos) -- (bl);
  \draw[arrow] (pos) -- (sv);
  \draw[arrow] (bl) |- (sv);
  \draw[arrow] (sv) -- (dow);
  \draw[arrow] (sv) -- (cmp);
  \draw[arrow] (cmp) -- (pbi);
\end{tikzpicture}%
}
\caption{Fluxo interno do \texttt{main.py} (simplificado): ETL alimenta o
solver; comparativo baseline e gravações seguem o resultado; DOW é condicional à
granularidade semanal; PBI consolidado depende da comparação produção--alocação.}
\label{fig:fluxo-main}
\end{figure}

\subsection{Casos de execução (A, B e parcial C)}
\label{sec:casos-execucao}

Os três modos abaixo correspondem ao que o manual recomenda operacionalmente.
O diagrama esquerdo é o fluxo completo; o central preserva edições manuais nos
CSVs; o direito é para manutenção pontual de uma fonte.

\begin{figure}[htbp]
\centering
\begin{tikzpicture}[
  node distance=4mm,
  casebox/.style={rectangle, draw, rounded corners=3pt, minimum width=4.6cm,
    text width=4.4cm, align=left, font=\small, inner sep=6pt},
  title/.style={font=\bfseries\small},
  arrow/.style={-{Stealth[length=2mm]}, thick}
]
  \node[casebox] (a) {\textbf{Fluxo A --- completo}\\[2mm]
    \texttt{./executar\_pipeline.sh}\\[1mm]
    {\footnotesize
    Regenera passos 1--6 a partir das bases brutas e roda
    \texttt{main.py}. Use quando PRIC, faturamento, produção ou cadastros mudaram.}};
  \node[casebox, right=8mm of a] (b) {\textbf{Fluxo B --- CSVs editáveis}\\[2mm]
    \texttt{./executar\_pipeline.sh {-}{-}sem-preparacao}\\[1mm]
    {\footnotesize
    Pula 1--6; usa os \texttt{inputs/*.csv} já na pasta (incl.\ ajustes manuais
    de negócio).}};
  \node[casebox, right=8mm of b] (c) {\textbf{Fluxo C --- parcial}\\[2mm]
    {\footnotesize
    Reexecutar só o(s) gerador(es) afetado(s) (ex.: novo período de pedidos) e,
    em seguida, \texttt{python3 main.py} ou o shell com
    \texttt{{-}{-}sem-preparacao}.}};
  \node[below=3mm of a, font=\footnotesize\itshape] (la) {$\rightarrow$ pasta \texttt{resultados/<ts>/}};
  \node[below=3mm of b, font=\footnotesize\itshape] (lb) {$\rightarrow$ idem};
  \node[below=3mm of c, font=\footnotesize\itshape] (lc) {$\rightarrow$ idem};
\end{tikzpicture}
\caption{Casos de execução: A regenera tudo; B respeita governança manual nos
CSVs; C exige disciplina para não misturar períodos inconsistentes entre
fontes.}
\label{fig:fluxo-casos}
\end{figure}

\textbf{Conversas entre componentes (resumo):}
\begin{itemize}[leftmargin=1.2cm]
    \item \textbf{config.yaml $\leftrightarrow$ todos os scripts:} período,
    estabelecimento, \texttt{paths} e regras do \texttt{modelo}; sem alinhamento,
    ETL e preparação geram insumos de ``mundos'' diferentes.
    \item \textbf{Preparação $\rightarrow$ ETL:} arquivos em \texttt{inputs/} são
    a camada que o \texttt{ETLPipeline} espera após a preparação (ou após edição
    manual no Fluxo B).
    \item \textbf{ETL $\rightarrow$ Otimizador:} contrato principal é
    \texttt{base\_otimizacao} + dicionários/séries de produção, excedente e
    pedidos garantidos.
    \item \textbf{Otimizador $\rightarrow$ Outputs:} resultado tabular alimenta
    baseline, arquivos de rodada, comparador externo e (em cadeia) consolidado
    PBI.
\end{itemize}

\section{Regras de negócio: decisão, justificativa e objetivo}

\subsection{Elegibilidade de SKU}

\textbf{Decisão adotada:} entram na otimização apenas SKUs ativos, no
estabelecimento alvo, com classe e embalagem válidas, além de preço, custo e
margem positivos.

\textbf{Justificativa:} permitir entrada irrestrita parece ampliar cobertura, mas
na prática contamina o problema com itens inexequíveis e degrada a leitura de
resultado.

\textbf{Objetivo:} preservar integridade econômica e operacional da recomendação.

\subsection{Pedidos primeiro, otimização depois}

\textbf{Decisão adotada:} com \texttt{atender\_pedidos=true}, pedidos
confirmados são reservados antes de otimizar.

\textbf{Justificativa:} compromisso comercial já assumido não pode competir em
pé de igualdade com volume discricionário de margem.

\textbf{Objetivo:} proteger nível de serviço e evitar que ganho financeiro
teórico comprometa atendimento real.

\textbf{Formulação operacional por classe:}
\[
E_c = \max(0, P_c - R_c)
\]
em que:
\begin{itemize}[leftmargin=1.2cm]
    \item \(P_c\): produção da classe no período;
    \item \(R_c\): volume reservado para pedidos da classe;
    \item \(E_c\): excedente disponível para otimização.
\end{itemize}

Com \texttt{capar\_reserva\_na\_producao=true}, o sistema impede que a reserva
ultrapasse a produção da classe, reforçando consistência física.

\subsection{Demanda histórica como freio de realismo}

\textbf{Decisão adotada:} o modelo pode limitar alocação por SKU com base
histórica (\texttt{media}, \texttt{maximo}, \texttt{percentil}) e fator de
ajuste.

\textbf{Justificativa:} sem esse limite, a otimização pode concentrar volume em
itens de alta margem sem absorção comercial.

\textbf{Objetivo:} equilibrar rentabilidade e viabilidade de sell-out.

\subsection{Fallbacks com transparência obrigatória}

\textbf{Decisão adotada:} quando faltam dados, o processo usa fallback
determinístico:
\begin{itemize}[leftmargin=1.2cm]
    \item preço: direto \(\rightarrow\) por item \(\rightarrow\) médio de classe \(\rightarrow\) médio geral;
    \item custo: direto \(\rightarrow\) médio de classe \(\rightarrow\) fallback parametrizado;
    \item embalagem: regex \(\rightarrow\) \texttt{embalagens\_override.xlsx}.
\end{itemize}

\textbf{Justificativa:} excluir todos os casos incompletos reduz cobertura e
limita valor prático.

\textbf{Objetivo:} ampliar aplicabilidade sem perder rastreabilidade.

\textbf{Garantia analítica:} colunas \texttt{origem\_preco},
\texttt{origem\_custo} e \texttt{origem\_embalagem} tornam explícita a
proveniência do dado usado.

\subsection{Comparador e PBI para leitura de causa}

\textbf{Decisão adotada:} outputs separam alocação, reserva e não atendimento, e
mostram produção bruta, estorno EAC e líquida.

\textbf{Justificativa:} quando essas camadas são misturadas, o usuário interpreta
``erro'' onde existe apenas diferença semântica.

\textbf{Objetivo:} reduzir ruído analítico e permitir leitura causal do resultado.

\subsection{SKUs fora da otimização com motivo explícito}

\textbf{Decisão adotada:} classificar ``fora da otimização'' por critério objetivo
(\texttt{quantidade\_alocada} \(\le 0\)) e exigir motivo principal.

\textbf{Justificativa:} classificações por rótulo textual geram falso positivo em
cenários agregados.

\textbf{Objetivo:} transformar exclusão em diagnóstico acionável.

\section{Decisões de modelagem e escolhas de insumos (síntese histórica)}

Esta seção consolida decisões recorrentes observadas no código e no histórico de
evolução do projeto. A intenção é registrar \textbf{o que foi decidido},
\textbf{por que foi decidido} e \textbf{qual efeito prático a decisão produz}.

\subsection{Evolução da modelagem (núcleo decisório)}

\textbf{1) Reserva tratada como prioridade estrutural, não como exceção}\\
Ao longo da evolução, a regra de pedidos migrou de tratamento auxiliar para
elemento estrutural do modelo. A classe passou a consumir reserva antes da
otimização do excedente, evitando ganho artificial de margem em cima de volume
já comprometido.

\textbf{2) Capacidade por classe com semântica física consistente}\\
As correções sucessivas nas restrições consolidaram a ideia de que a capacidade
válida para otimizar é a capacidade líquida após reservas. Essa escolha reduz
situações em que a solução matemática era ótima, mas fisicamente inexequível.

\textbf{3) Demanda histórica como mecanismo de realismo comercial}\\
A inclusão e refinamento dos métodos de demanda (\texttt{maximo}, \texttt{media},
\texttt{percentil}, com fator) mostram a decisão de evitar concentração irreal em
SKUs de alta margem sem respaldo de mercado.

\textbf{4) Baseline como instrumento de comparação justa}\\
A evolução do comparativo baseline eliminou vieses importantes (como média simples
indevida ou mistura com linhas de reserva), para que o ganho informado reflita
realocação de mix e não efeito de metodologia.

\textbf{5) Granularidade D/S/M como decisão operacional}\\
O suporte a diário, semanal e mensal não é apenas conveniência técnica; é uma
decisão de aderência ao ciclo real de planejamento. O modelo mantém o mesmo
núcleo matemático, mudando a referência temporal dos insumos e relatórios.

\subsection{Decisões sobre inputs brutos e derivados}

\textbf{1) Custo bruto (fonte primária)}\\
O custo é apurado da base bruta de custos PRIC (Parquet), com filtros de
estabelecimento, período e qualidade mínima de dado, gerando o insumo editável
\texttt{inputs/custos\_sku.csv}.

\textbf{2) Preço operacional (decisão vigente de negócio)}\\
Pela regra atualmente adotada no projeto, o preço operacional usado no pipeline
é reconstruído via processo de extração que hoje está acoplado ao mesmo insumo
bruto referenciado em \texttt{paths.custos}, gerando
\texttt{inputs/precos\_sku\_embalagem.csv}.  
Essa escolha foi mantida para contornar inconsistências identificadas na base de
faturamento para formação de preço.

\textbf{3) Faturamento bruto permanece crítico para outras frentes}\\
Mesmo quando não é a base primária de preço operacional, o faturamento bruto
segue central para:
\begin{itemize}[leftmargin=1.2cm]
    \item demanda histórica;
    \item compatibilidade SKU-embalagem;
    \item enriquecimento descritivo e rastreabilidade em outputs.
\end{itemize}

\textbf{4) CSVs derivados como contrato operacional}\\
A adoção de \texttt{precos\_sku\_embalagem.csv}, \texttt{custos\_sku.csv},
\texttt{demanda\_historica.csv} e \texttt{pedidos\_clientes.csv} como camada
editável explicita uma decisão metodológica: separar cálculo bruto de governança
de negócio, permitindo ajustes controlados sem alterar código.

\subsection{Decisões de transparência e auditoria}

\textbf{1) Rastreabilidade de origem (preço/custo/embalagem)}\\
A propagação de colunas de origem em toda a cadeia (ETL, resultado, comparador e
PBI) foi uma decisão para reduzir discussões subjetivas sobre ``de onde veio o
número''.

\textbf{2) Produção líquida como base canônica de comparação}\\
A comparação passou a privilegiar produção líquida (bruta - EAC), tornando
coerente a leitura de diferença entre produção e alocação.

\textbf{3) ``Fora da otimização'' por critério objetivo}\\
A regra consolidada usa \(\texttt{quantidade\_alocada} \le 0\), com motivo
explícito. Isso evita falso-positivo causado por rótulos compostos e melhora a
ação corretiva.

\textbf{4) Saída por rodada e base física persistida}\\
Salvar cada execução em diretório timestampado e persistir
\texttt{base\_otimizacao} em \texttt{parquet/csv/xlsx} foi decisão de
reprodutibilidade: toda conclusão pode ser reconstituída.

\subsection{Leitura executiva das decisões}

Em conjunto, essas decisões apontam um padrão claro:
\begin{itemize}[leftmargin=1.2cm]
    \item o projeto prioriza \textbf{consistência operacional} sobre atalho estatístico;
    \item prioriza \textbf{explicabilidade} sobre ``caixa-preta'' de resultado;
    \item prioriza \textbf{comparação justa} sobre ganho inflado por metodologia;
    \item prioriza \textbf{governança de dados} por meio de insumos derivados auditáveis.
\end{itemize}

\section{Parametrização em \texttt{config.yaml}: negócio e implicações}
\label{sec:param-config}

O \texttt{config.yaml} é o \textbf{contrato único} entre gestão, dados e código:
alterá-lo muda o \textit{mundo} em que o modelo decide (período, unidade
geográfica, prioridade comercial e limites de realismo). Os blocos abaixo
espelham a estrutura do arquivo; parâmetros puramente numéricos do solver são
citados de forma resumida ao final.

\subsection{Bloco \texttt{dados} --- janela temporal, unidade e qualidade}

\textbf{Ano/mês e janelas (\texttt{ano\_custo}, \texttt{mes\_custo},
\texttt{meses\_janela\_custo}, espelhados para preço):} definem o intervalo
sobre o qual custos, preços médios e (em parte) demanda histórica são
calculados. \textbf{Implicação:} janela curta reage a mudanças recentes, mas é
mais ruidosa; janela longa estabiliza números, mas pode ``congelar'' estrutura
de custo ou preço já superada.

\textbf{Estabelecimentos (\texttt{estab\_custo}, \texttt{estab\_preco},
\texttt{estabelecimentos}):} amarram custo, preço e filtros de faturamento ao
mesmo conjunto de unidades. \textbf{Implicação:} trocar estabelecimento sem
revisar insumos de produção e pedidos pode gerar otimização para um hub que não
corresponde ao cenário operado.

\textbf{Preço --- outliers (\texttt{remover\_outliers\_preco},
\texttt{percentil\_inferior\_preco}, \texttt{percentil\_superior\_preco}):}
recortam caudas extremas antes da média. \textbf{Implicação:} reduz distorção
por erros ou vendas atípicas; porém percentis muito apertados podem ``achatar''
realidade de itens estratégicos com poucos lançamentos.

\textbf{Referência de produção (\texttt{semana\_ref} quando
\texttt{granularidade\_demanda} = S; \texttt{data\_ref} quando D ou M):} fixa o
período cujo volume produtivo alimenta capacidade por classe. \textbf{Implicação:}
deve estar alinhado ao planejamento (semana/dia/mês que a gestão trata como
``a produção a distribuir'').

\textbf{Referência de pedidos (\texttt{pedidos\_semana\_ref},
\texttt{pedidos\_data\_ref}):} são \textbf{independentes} da referência de
produção: permitem, por exemplo, produção de uma semana e carteira de pedidos
de outra. \textbf{Implicação:} útil quando ciclos comercial e fabril não
coincidem; exige disciplina explícita para não comparar inadvertidamente
períodos incongruentes na leitura de resultados.

\textbf{Perfil DOW (\texttt{dow\_ano\_ref}, \texttt{dow\_mes\_ref},
\texttt{dow\_meses\_janela}):} usados na granularidade \textbf{semanal} para
distribuir demanda histórica pelo dia da semana. \textbf{Implicação:} se
omitidos, o código recai na janela de custo; ajuste quando o padrão semanal
relevante for outro (sazonalidade, feriados).

\textbf{Filtros auxiliares (\texttt{sku\_minimo\_transacoes},
\texttt{tipo\_estoque}):} elevam o piso de evidência estatística e o recorte de
estoque. \textbf{Implicação:} valores altos de transações mínimas reduzem ruído,
mas podem excluir SKUs de nicho ainda estratégicos.

\subsection{Bloco \texttt{modelo} --- regras da otimização}

\textbf{Pedidos e excedente (\texttt{atender\_pedidos},
\texttt{usar\_apenas\_excedente}):} com ambos verdadeiros, reserva-se primeiro o
comprometido e só o saldo entra no solver. \textbf{Implicação:} aderência a
nível de serviço; desligar \texttt{atender\_pedidos} transforma o problema em
``puro mix'' (aceitável só em estudos, não em operação com carteira firme).

\textbf{\texttt{capar\_reserva\_na\_producao}:} impede reserva acima da produção
da classe. \textbf{Implicação:} coerência física e priorização explícita quando
há competição entre pedidos dentro da mesma classe (ver formulação na seção de
regras de negócio).

\textbf{Demanda histórica (\texttt{considerar\_demanda\_historica},
\texttt{tipo\_calculo\_demanda}, \texttt{fator\_demanda\_maxima},
\texttt{percentil\_demanda}):} quando ativo, limita quanto cada SKU pode receber
em função do histórico (\texttt{media}, \texttt{maximo} ou
\texttt{percentil}), multiplicado pelo fator. \textbf{Implicações:}
\texttt{maximo} é conservador para sell-out; \texttt{media} permite mais
flexibilidade; \texttt{percentil} calibra risco (percentil alto = teto mais
permissivo). O fator \(>1\) amplia deliberadamente o teto (crescimento esperado);
fator \(=1\) ancora no histórico ``cru''.

\textbf{\texttt{granularidade\_demanda} (S, M, D):} define a periodicidade dos
insumos e relatórios. \textbf{Implicação:} deve refletir o ritmo de decisão
(semana operacional vs. fechamento mensal vs. dia); misturar mentalmente
granularidades na interpretação dos gráficos gera os ``percentuais estranhos''
mencionados na seção de erros recorrentes.

\textbf{\texttt{tipo\_objetivo}:} \texttt{maximizar\_margem} (padrão) ou
\texttt{minimizar\_custos}. \textbf{Implicação:} a primeira alinha com rentabilidade
unitária; a segunda pode fazer sentido em exercícios de eficiência de custo,
desalinhando a narrativa comercial se usada sem consenso.

\textbf{\texttt{variaveis\_continuas}:} volumes fracionários vs.\ inteiros.
\textbf{Implicação:} contínuo acelera e suaviza soluções; inteiro aproxima
unidades físicas (caixas), com custo de tempo de solver maior.

\subsection{Bloco \texttt{negocio}}

\textbf{\texttt{filtrar\_granjas}:} lista de estabelecimentos (identificadores
como na base) cuja demanda histórica entra no cálculo de limites. \textbf{Implicação:}
deve ser coerente com a estratégia de mercado (incluir só a rede relevante evita
que filiais fora do escopo ``puxem'' ou diluam o perfil de demanda).

\subsection{Bloco \texttt{paths} --- onde o pipeline lê e grava}

Cada chave aponta para planilhas/Parquet/CSV de entrada e pastas de saída
(\texttt{logs\_dir}, \texttt{output\_dir}). \textbf{Implicação:} mudar caminho
sem atualizar o arquivo físico quebra a rodada; manter \texttt{pedidos},
\texttt{producao}, \texttt{classes} e derivados (\texttt{precos},
\texttt{custos}) consistentes com o período configurado em \texttt{dados} é
requisito de governança, não detalhe técnico.

\subsection{Bloco \texttt{fallbacks}}

Valores usados quando o ETL não encontra custo válido ou precisa de referência
de embalagem. \textbf{Implicação:} alto uso de fallback nas colunas
\texttt{origem\_*} sinaliza dado frágil; negócio deve tratar como alerta, não
como ``preço/custo oficial''.

\subsection{Bloco \texttt{solver} (resumo)}

Limites de tempo (\texttt{time\_limit\_ms}), \texttt{gap\_percentage} e threads
afetam \textbf{quão longo} e \textbf{quão próximo do ótimo} o programa para.
\textbf{Implicação} para gestão: em bases muito grandes, relaxar tempo ou aceitar
gap um pouco maior pode ser preferível a interrupção sem solução utilizável.

\subsection{Síntese: o que revisar antes de cada rodada}

Para decisão de negócio, a ordem prática de revisão é:
\begin{enumerate}[leftmargin=1.2cm]
    \item \texttt{dados}: estabelecimento(s), janelas de custo/preço e referências
    de produção \emph{e} pedidos;
    \item \texttt{modelo}: tríade pedidos/excedente/cap, demanda histórica + fator,
    granularidade e objetivo;
    \item \texttt{negocio.filtrar\_granjas} alinhado ao escopo comercial;
    \item \texttt{paths} coerentes com os arquivos que a equipe pretende usar.
\end{enumerate}

\subsection{Matriz parametro -> impacto esperado no output}

\begin{longtable}{>{\raggedright\arraybackslash}p{5.2cm}>{\raggedright\arraybackslash}p{5.0cm}>{\raggedright\arraybackslash}p{5.0cm}}
\toprule
\textbf{Parametro} & \textbf{Efeito direto} & \textbf{Onde aparece} \\
\midrule
\endhead
\url{modelo.granularidade_demanda} & Troca o recorte temporal (D/S/M) e a referencia de periodo. & Nome da pasta (\texttt{\_D}, \texttt{\_S} ou \texttt{\_M}), comparador, parametros do baseline. \\
\url{dados.data_ref} / \url{dados.semana_ref} & Define o periodo de producao usado no comparador e rastreabilidade. & \url{comparacao_producao_alocacao}, \url{resumo_por_classe}. \\
\url{dados.pedidos_data_ref} / \url{dados.pedidos_semana_ref} & Define o periodo da carteira de pedidos (independente da producao). & \url{quantidade_reservada}, \url{quantidade_nao_atendida_pedido}. \\
\url{modelo.atender_pedidos} & Liga/desliga reserva previa de pedidos. & Resultado completo, comparador, aba de SKUs fora. \\
\url{modelo.capar_reserva_na_producao} & Impede reservar acima da producao da classe. & Deficit de pedido e priorizacao em reservas. \\
\url{modelo.considerar_demanda_historica} + \texttt{fator} & Ajusta teto de alocacao por SKU e baseline com cap. & \url{limite_demanda_historica}, \url{volume_baseline}. \\
\url{dados.estabelecimentos} & Define universo de dados e elegibilidade por estabelecimento. & Todos os outputs + metadados de estabelecimento. \\
\bottomrule
\end{longtable}

\section{Como executar: método operacional}
\label{sec:como-executar}

\subsection{Três fluxos de execução}

\textbf{Fluxo A -- Completo (recomendado quando bases brutas mudam):}
\begin{verbatim}
source venv/bin/activate
./executar_pipeline.sh
\end{verbatim}

\textbf{Fluxo B -- Preservar CSVs editáveis (quando houve ajuste manual):}
\begin{verbatim}
source venv/bin/activate
./executar_pipeline.sh --sem-preparacao
\end{verbatim}

\textbf{Fluxo C -- Parcial (quando só uma fonte mudou):}
reexecute apenas os geradores afetados e, em seguida, rode \texttt{main.py}.

\subsection{Ritual mínimo antes de rodar}

\begin{enumerate}[leftmargin=1.2cm]
    \item validar período, estabelecimento, referências de produção/pedidos e bloco
    \texttt{modelo} em \texttt{config.yaml} (ver seção dedicada);
    \item confirmar se deve preservar ou recalcular CSVs de \texttt{inputs/};
    \item checar disponibilidade de dados para o período escolhido;
    \item executar o fluxo adequado (A/B/C);
    \item revisar outputs na pasta da rodada.
\end{enumerate}

\section{Leitura dos resultados: como argumentar com segurança}
\label{sec:leitura-resultados}

\subsection{Resultado detalhado}

O \texttt{resultado\_realocacao\_completo} mostra o ``o que foi decidido'' no
nível item/embalagem.

\subsection{Ganho versus baseline}

A \texttt{auditoria\_baseline} responde se a realocação efetivamente criou valor,
comparando o cenário otimizado com o cenário de referência.

\textbf{Leitura obrigatoria das colunas de volume no baseline:}
\begin{itemize}[leftmargin=1.2cm]
    \item \texttt{volume\_baseline\_antes\_cap}: volume de referencia sem aplicar limite de demanda;
    \item \texttt{volume\_baseline}: volume apos cap + redistribuicao;
    \item \texttt{volume\_redistribuido} \(=\) \texttt{volume\_baseline} \(-\) \texttt{volume\_baseline\_antes\_cap}.
\end{itemize}

\textbf{Exemplo numerico simples (em ovos):}
\begin{itemize}[leftmargin=1.2cm]
    \item SKU A: antes do cap = 10.000, limite = 8.000, depois do cap = 8.000, redistribuido = -2.000;
    \item SKU B: antes do cap = 6.000, limite = 9.000, recebe parte do excedente, depois = 7.500, redistribuido = +1.500.
\end{itemize}
Interprete valores negativos como \textbf{SKU capado} e positivos como \textbf{SKU que recebeu redistribuicao}.

\subsection{Produção versus alocação}

O \texttt{comparacao\_producao\_alocacao} é a principal fonte para discussão com
operação e comercial. As colunas-chave são:
\begin{itemize}[leftmargin=1.2cm]
    \item \texttt{quantidade\_alocada},
    \item \texttt{quantidade\_reservada},
    \item \texttt{quantidade\_nao\_atendida\_pedido},
    \item \texttt{quantidade\_produzida\_bruta},
    \item \texttt{quantidade\_estorno\_eac},
    \item \texttt{quantidade\_produzida\_liquida\_modelo}.
\end{itemize}

\subsection{Fora da otimização}

O \texttt{skus\_fora\_otimizacao} informa quem ficou sem alocação e por qual
motivo, permitindo ação corretiva dirigida.

\textbf{Dicionario pratico de motivos (mais comuns):}
\begin{itemize}[leftmargin=1.2cm]
    \item \texttt{bloqueado\_por\_regra\_tipo\_granel}: SKU fora por regra de tipo no cadastro;
    \item \texttt{bloqueado\_por\_regra\_tipo\_exportacao}: SKU fora por tipo exportacao;
    \item \texttt{bloqueado\_por\_regra\_status\_inativo}: SKU nao elegivel no estabelecimento;
    \item \texttt{bloqueado\_por\_regra\_estab\_fora\_escopo}: item fora dos estabelecimentos configurados;
    \item \texttt{priorizacao\_margem\_intraclasse}: SKU perdeu volume para outro SKU da mesma classe com margem superior;
    \item \texttt{sem\_producao\_no\_periodo}: sem producao liquida para alocar no periodo.
\end{itemize}

\textbf{Acao recomendada por motivo:}
\begin{itemize}[leftmargin=1.2cm]
    \item motivos de regra: revisar cadastro/carteira, nao o solver;
    \item priorizacao de margem: avaliar estrategia comercial da classe;
    \item sem producao: tratar planejamento/producao, nao parametrizacao de demanda.
\end{itemize}

\subsection{Consumo analítico}

O \texttt{pbi\_consolidado} organiza dados para BI sem perda de contexto de
parâmetros e rastreabilidade.

\section{Rastreabilidade e reprodutibilidade}
\label{sec:rastreabilidade}

Cada execução cria uma pasta própria de rodada:
\[
\texttt{resultados/<timestamp\_rodada>\_<D|S|M>/}
\]

Além disso, a base efetivamente usada na otimização é persistida em:
\begin{itemize}[leftmargin=1.2cm]
    \item \texttt{base\_otimizacao\_<timestamp>.parquet}
    \item \texttt{base\_otimizacao\_<timestamp>.csv}
    \item \texttt{base\_otimizacao\_<timestamp>.xlsx}
\end{itemize}

Esse desenho não é detalhe técnico: é a condição para reproduzir uma decisão e
explicar, de forma verificável, o que entrou no solver.

\section{Erros recorrentes e interpretação correta}

\subsection*{``Não há registros para o período''}
Normalmente indica incompatibilidade entre período configurado e dados disponíveis
nos inputs (ou ausência de dados no estabelecimento selecionado).

\subsection*{Percentual de alocação aparentemente alto}
Costuma ocorrer quando há desalinhamento entre granularidade/periodização dos
insumos usados no ETL e aqueles usados no comparador.

\subsection*{Muitos SKUs com \texttt{SEM\_EMBALAGEM}}
Indica baixa captura automática por regex; deve-se complementar
\texttt{embalagens\_override.xlsx}.

\section{Checklist executivo da rodada}
\label{sec:checklist-executivo}

\begin{itemize}[leftmargin=1.2cm]
    \item \texttt{config.yaml}: período, estabelecimento, referências produção/pedidos
    e regras de demanda/objetivo conferidos;
    \item pedidos filtrados na granularidade correta;
    \item produção líquida coerente (bruta - estorno);
    \item origem de dados visível (preço, custo, embalagem);
    \item SKUs fora da otimização com motivo explícito;
    \item ganho versus baseline consistente com limites de demanda.
\end{itemize}

\textbf{Criterios objetivos de aceite (go/no-go):}
\begin{itemize}[leftmargin=1.2cm]
    \item classes com producao liquida relevante devem fechar \texttt{alocacao + reserva} proximo da producao liquida;
    \item itens com \texttt{fora\_otimizacao\_sem\_motivo\_claro} devem ser excecao e investigados;
    \item mudanca de periodo/estabelecimento no \texttt{config.yaml} exige coerencia com os CSVs preservados;
    \item leitura de ganho deve usar baseline apos cap (\texttt{volume\_baseline}), nao apenas baseline antes do cap.
\end{itemize}

\section{Conclusão}

Este modelo deve ser entendido como um sistema de governança decisória, e não
apenas como um cálculo de alocação. Ele organiza conflito entre margem,
capacidade, demanda e atendimento em um processo auditável.

Em termos práticos, seu valor não está apenas em ``quanto alocou'', mas em
\textbf{por que alocou}, \textbf{o que deixou de alocar}, \textbf{quais limites
respeitou} e \textbf{qual resultado financeiro produziu frente ao baseline}.

Quando essa lógica é seguida com disciplina de dados e leitura adequada dos
outputs, a empresa ganha previsibilidade, consistência e qualidade de decisão.

\end{document}


